\chapter{Annexes}

\section{Arborescence du projet}
\begin{lstlisting}[caption={Arborescence simplifiée},label={lst:arbre}]
jee-ecommerce-jsf-jpa/
  pom.xml, README.md, .gitignore
  db/create-database.sql
  docs/diagramme-classes.md, docs/rapport/
  src/main/java/ma/fstt/
    entities/  exception/  util/  persistence/  service/  jpa/   (atelier 2)
    web/       15 classes : filtres, beans, UserSession, Affichage
  src/main/resources/
    META-INF/persistence.xml, META-INF/beans.xml, db.properties.example
  src/main/webapp/
    WEB-INF/web.xml, faces-config.xml, templates/layout.xhtml
    resources/components/  csrf, icone, visuel, vide, champ
    *.xhtml (8 pages) et admin/*.xhtml (4 pages)
    assets/css/app.css, assets/js/app.js
  src/test/java/ma/fstt/  Verifications + 7 lots de contrôles
\end{lstlisting}

\section{Comptes de démonstration}
\begin{table}[htbp]
\centering
\caption{Comptes créés par \texttt{SchemaTool}}\label{tab:comptes}
\begin{tabular}{@{}lll@{}}
\toprule
\textbf{Profil} & \textbf{Adresse} & \textbf{Mot de passe}\\
\midrule
Administrateur & \code{admin@example.com} & \code{Admin\#2026}\\
Client & \code{yasmine@example.com} & \code{Client\#2026}\\
\bottomrule
\end{tabular}
\end{table}
Un second client, \code{omar@example.com}, sert au scénario de refus d'accès. Ces comptes n'existent que dans la base
locale de démonstration.

\section{Diagrammes Mermaid}
Le fichier \code{docs/diagramme-classes.md} contient deux diagrammes que GitHub affiche directement. Le premier est le modèle
du domaine ; le second la couche JSF.

\fig{01-diagramme-metier}{Diagramme Mermaid du modèle métier}{fig:mermaid-metier}
\fig{02-diagramme-beans}{Diagramme Mermaid de la couche JSF}{fig:mermaid-beans}
